\section{Alur kerja Petugas/PIC di mobile}

Bagian ini mengikuti alur tugas SOP dari beranda sampai tugas ditandai selesai. Pada layar Petugas, laporan aktivitas dan penyelesaian tugas adalah dua tindakan terpisah.

\subsection{Periksa beranda, tugas perlu tindakan, dan insiden}

\begin{enumerate}
  \item Setelah masuk, buka \textbf{Beranda}. Bagian atas menampilkan cuaca serta ringkasan subblok pra-tanam, penanaman, dan panen.
  \item Periksa kartu \textbf{Insiden aktif}. Insiden yang belum ditindaklanjuti dapat berasal dari laporan Petugas pada lahan yang sama; buka detail atau pilih \textbf{Lihat semua} untuk meninjau informasi dan tindak lanjut.
  \item Pada bagian \textbf{Tugas perlu tindakan} (Gambar~\ref{fig:petugas-beranda}), pilih kartu tugas untuk membuka detailnya, atau pilih \textbf{Lihat semua} untuk masuk ke daftar lengkap tugas. Indikator progres (\emph{X/Y selesai}), tanggal jadwal, target HST, dan lencana status membantu membedakan pekerjaan yang harus diprioritaskan.
  \item Jika tugas atau insiden yang diharapkan tidak terlihat, pastikan lahan yang ditugaskan benar dan periksa menu Tugas bersama Admin.
\end{enumerate}

\guidemobileimage[0.40\linewidth]{GambarPetugas1.png}{Beranda Petugas: tinjau insiden aktif dan pilih tugas pada daftar tugas. Catatan: pada versi terbaru, judul daftar To-do hari ini telah diperbarui menjadi Tugas perlu tindakan disertai tautan Lihat semua dan tanggal jadwal pada kartu tugas.}{fig:petugas-beranda}

\subsection{Buka dan pahami detail tugas}

\begin{enumerate}
  \item Pilih tugas dari daftar \textbf{Tugas perlu tindakan} di Beranda atau melalui tab \textbf{Tugas}.
  \item Periksa nama tugas, lahan/blok/subblok, \textbf{Jadwal tugas}, \textbf{Target tugas} (HST), \textbf{Varietas}, \textbf{Estimasi HOK}, serta \textbf{Jendela pelaporan}.
  \item Baca \textbf{Instruksi SOP} sebelum mulai bekerja. Instruksi menjelaskan cara kerja atau hal yang harus diperiksa.
  \item Jika pencatatan tersedia (tanggal tanam aktual sudah tercatat dan tugas belum selesai/ditutup), tekan \textbf{Catat aktivitas tugas ini} untuk membuat laporan yang terhubung dengan tugas.
\end{enumerate}

\guideimage{GambarPetugas2.png}{Baca target HST dan instruksi SOP sebelum memulai pekerjaan, lalu buka form dari tombol Catat aktivitas tugas ini.}{fig:petugas-detail-tugas}

\guideimage{GambarPetugas3.png}{Contoh tampilan detail tugas dan instruksi yang menjadi acuan pelaksanaan.}{fig:petugas-detail-instruksi}

\guideimage{GambarPetugas7.png}{Tindakan untuk memulai pencatatan laporan pada tugas SOP melalui tombol Catat aktivitas tugas ini.}{fig:petugas-tombol-catat}

\subsection{Isi form aktivitas yang terhubung ke tugas}

\begin{enumerate}
  \item Pastikan ringkasan tugas yang terhubung sudah benar. \textbf{Tanggal kerja} harus sesuai hari pelaksanaan dan berada dalam jendela pelaporan yang ditampilkan.
  \item Isi \textbf{Jumlah pekerja} serta \textbf{Luas area} dan satuannya sesuai pekerjaan. Kolom dapat berbeda menurut kategori tugas (misalnya \textbf{Jumlah pekerja} tidak ditampilkan pada pengamatan perkembangan tanaman); isi kolom yang tampil dan wajib.
  \item Lampirkan minimal satu \textbf{foto bukti pekerjaan} melalui \textbf{Tambah foto}. Form menerima JPG, PNG, atau WebP, paling banyak lima foto, dengan batas ukuran 15~MB per foto (ukuran foto disesuaikan otomatis sebelum diunggah).
  \item \textbf{Tambah catatan} dapat diisi jika diperlukan. Periksa kembali data dan bukti, lalu tekan \textbf{Simpan aktivitas}.
  \item Jika unggahan atau penyimpanan gagal, ikuti pesan pada layar dan kirim ulang setelah memperbaiki isian atau koneksi.
\end{enumerate}

\guideimage{GambarPetugas4.png}{Form aktivitas tugas. Kolom menyesuaikan jenis pekerjaan; foto bukti wajib dilampirkan sebelum laporan dikirim.}{fig:petugas-form-aktivitas}

\guidemobileimage[0.44\linewidth]{GambarPetugas8.png}{Contoh pengisian laporan untuk tugas SOP sebelum menekan Simpan aktivitas.}{fig:petugas-form-terisi}

\subsection{Tinjau laporan yang tersimpan dan selesaikan tugas}

\begin{enumerate}[series=petugas-selesai]
  \item Setelah penyimpanan berhasil, periksa detail tugas atau buka tab \textbf{Laporan} untuk memastikan aktivitas muncul. Laporan tugas menampilkan relasi ke tugas dan ringkasan hasil yang dikirim.
  \item Tersimpannya laporan \textbf{tidak otomatis menyelesaikan tugas}. Status tugas berubah menjadi \textbf{Dikerjakan} dan dapat tetap demikian selama masih ada pekerjaan yang harus dituntaskan (Gambar~\ref{fig:petugas-status-dikerjakan}).
\end{enumerate}

\guidemobileimage[0.46\linewidth]{GambarPetugas10.png}{Setelah laporan terkait dikirim, tugas tampak berstatus Dikerjakan; status ini belum berarti tugas selesai.}{fig:petugas-status-dikerjakan}

\begin{enumerate}[resume=petugas-selesai]
  \item Setelah seluruh pekerjaan benar-benar selesai, buka kembali detail tugas, isi \textbf{Catatan penyelesaian (opsional)} bila perlu, lalu tekan \textbf{Selesaikan tugas} (Gambar~\ref{fig:petugas-selesaikan-tugas}).
  \item Tinjau kembali status setelah konfirmasi. Menandai selesai merupakan sinyal progres untuk Admin dan sistem; setelah tugas selesai, laporan baru tidak dapat ditambahkan pada tugas tersebut.
\end{enumerate}

\guideimage{GambarPetugas6.png}{Contoh halaman detail laporan yang tersimpan. Periksa informasi tugas terkait untuk memastikan laporan terhubung dengan pekerjaan SOP yang benar.}{fig:petugas-detail-laporan}

\guidemobileimage[0.42\linewidth]{GambarPetugas9.png}{Tugas berstatus Dikerjakan dengan laporan terkait. Tekan Selesaikan tugas setelah seluruh pekerjaan tuntas.}{fig:petugas-selesaikan-tugas}

\subsection{Lihat semua tugas dan gunakan filter}

\begin{enumerate}
  \item Buka tab \textbf{Tugas} untuk melihat daftar tugas lintas blok dan subblok pada lahan penugasan. Gunakan tab \textbf{Tugas berjalan} untuk pekerjaan aktif atau \textbf{Riwayat subblok} untuk meninjau subblok yang sudah selesai/dibatalkan.
  \item Gunakan ringkasan jumlah tugas \textbf{Terlambat}, \textbf{Hari ini}, dan \textbf{Mendatang} sebagai gambaran beban kerja.
  \item Cari nama tugas atau subblok, pilih status (\textbf{Perlu tindakan}, \textbf{Terlambat}, \textbf{Hari ini}, \textbf{Mendatang}, \textbf{Ditutup}, atau \textbf{Semua}), lalu tekan tombol \textbf{Filter} untuk menyaring berdasarkan kategori tugas, blok, atau subblok. Gunakan \textbf{Reset filter} dan pengatur \textbf{Per halaman} (\texttt{10}, \texttt{20}, atau \texttt{50}) bila daftar panjang.
  \item Buka kartu tugas untuk melihat instruksi, mengirim laporan, atau menandai penyelesaian (Gambar~\ref{fig:petugas-daftar-tugas}). Daftar dapat mencakup tugas yang sudah berstatus \textbf{Dikerjakan} agar Petugas bisa melanjutkan tindak lanjut.
\end{enumerate}

\guidemobileimage[0.34\linewidth]{GambarPetugas11.png}{Daftar tugas lengkap dengan pencarian, filter status, dan daftar lintas blok/subblok.}{fig:petugas-daftar-tugas}

\subsection{Lihat lahan dan catat aktivitas mandiri}

\begin{enumerate}
  \item Buka tab \textbf{Lahan} (\textbf{Lahan saya}), lalu pilih blok untuk melihat alamat/wilayah, luas, kapasitas blok, tampilan subblok \textbf{Berjalan} atau \textbf{Riwayat}, serta ringkasan progres. Bila diberi izin dan sisa luas blok masih tersedia, Petugas juga dapat menekan \textbf{Tambah subblok} dari halaman ini.
  \item Untuk pekerjaan di luar daftar tugas SOP, tekan tombol \textbf{+} pada \textbf{Beranda} atau \textbf{Laporan} untuk membuka dialog \textbf{Tambah aktivitas mandiri}. Memilih kategori \textbf{Penanaman} akan mengarahkan ke form pencatatan penanaman, sedangkan kategori lainnya membuka form aktivitas mandiri (isi tanggal/waktu, lokasi, jumlah pekerja opsional, foto bukti, dan rincian yang diminta lalu tekan \textbf{Simpan laporan}).
  \item Pastikan aktivitas yang baru disimpan muncul di riwayat tab \textbf{Laporan} (Gambar~\ref{fig:petugas-aktivitas-mandiri}).
  \item Pada tab \textbf{Laporan}, cari berdasarkan aktivitas atau subblok dan gunakan filter \textbf{Tanggal kerja}, \textbf{Kategori}, \textbf{Sumber laporan} (\textbf{Semua sumber}, \textbf{Dari SOP}, \textbf{Mandiri}), atau centang \textbf{Input saya}. Buka kartu laporan untuk melihat detailnya.
\end{enumerate}

\guideimage{GambarPetugas5.png}{Contoh form aktivitas mandiri di luar jadwal tugas SOP. Gunakan alur ini untuk pekerjaan yang memang tidak terkait tugas terjadwal.}{fig:petugas-aktivitas-mandiri}
